\section*{Everything in one table}

Four experiments, two of them exact and two estimated. ``Exact'' means
computed by enumerating posteriors --- no sampler, no fitting, no replicates,
so there is no error bar to quote. The sections give the detail.

\begin{center}
\footnotesize
\setlength{\tabcolsep}{4pt}
\begin{tabular}{@{}p{5.9cm}p{3.5cm}p{5.2cm}@{}}
\toprule
Claim & Evidence & Result \\
\midrule

\multicolumn{3}{@{}l}{\textbf{Exact (\S\ref{sec:channels}) --- computed, not estimated}}\\
\addlinespace[2pt]
Bijective $g$: observing $h$ adds nothing
 & 20 maps, enumerated
 & $\mathbf{0.000000}$ for $z^3$ and $z^5$; non-zero for every other map \\
Linear $g$: the direct model is already right
 & identity, sum
 & $\mathbf{0.000000}$ \\
Size of the information channel
 & 20 maps
 & $0$ (injective) to $0.14$ \\
Size of the non-information gap
 & 20 maps
 & $0.24$ to $0.43$, including where no information exists \\
\quad of which: hyperparameter estimation
 & identity, sum (linear)
 & $0.264$, $0.183$ --- no model shape possible there \\
\quad of which: model shape, net of that
 & 20 maps
 & $+0.076$ mean excess over the baseline \\
\textbf{Which term dominates}
 &
 & \textbf{hyperparameter estimation}; model shape ($0.076$) and information ($0.071$) are the same order \\
\addlinespace[3pt]

\multicolumn{3}{@{}l}{\textbf{Predictive accuracy (\S\ref{sec:exactposteriors}) --- CRPS, 240 fits per map}}\\
\addlinespace[2pt]
Composite predicts better than Direct
 & 20 maps $\times$ 240 fits
 & every map, $t \ge 3.6$; $207$--$240$ wins out of $240$ \\
Linear $g$ gains least
 & all four sample sizes
 & $+0.0050$ against $0.019$--$0.037$ elsewhere \\
\textbf{Crossing injective $\to$ non-injective}
 & difference-in-differences, even rows as control
 & $\mathbf{+0.0190 \pm 0.0080}$, $t=2.37$ \\
The effect is binary, not graded in $\bar n$
 & even-degree rows
 & flat as $\bar n$ climbs $2.0 \to 3.9$ \\
\addlinespace[3pt]

\multicolumn{3}{@{}l}{\textbf{Multi-objective (\S\ref{sec:mobo}) --- hypervolume, the paper's own solvers}}\\
\addlinespace[2pt]
The advantage transfers to hypervolume
 & 4 conditions $\times$ 32 runs
 & $+17\%$ to $+31\%$; $26$--$31$ wins of $32$ \\
Gain is large even when $g_1$ is linear
 & 32 runs
 & $+17.3\%$, $31/32$ --- where information is provably zero \\
Misspecification step ($z \to \sinh$)
 & paired
 & $+6.85 \pm 5.97$ \quad (n.s.) \\
Information step ($\sinh \to z^2$)
 & paired
 & $+1.63 \pm 6.63$ \quad (n.s.) \\
\textbf{Their ratio}
 &
 & $\approx 4{:}1$, matching the exact $3$--$7\times$ \\
\addlinespace[3pt]

\multicolumn{3}{@{}l}{\textbf{Structural results (\S\ref{sec:inseparable}) --- proved and measured}}\\
\addlinespace[2pt]
In one dimension, non-injective $\Rightarrow$ $g$ folds $\Rightarrow$ $y(x)$ gains turning points
 & counted on every map
 & injective maps: $1$ turn, as $h$ itself. Others: $2$--$3$ \\
An injective $g$ is monotone, so its optimum sits at an endpoint
 & --- & regret cannot compare the two classes fairly \\
\bottomrule
\end{tabular}
\end{center}

\paragraph{The one-sentence version.} Bijectivity decides, exactly, whether
observing the intermediates carries any information --- and that is the
smaller reasons the composite model wins. The largest single term is neither
information nor model shape but the fact that the composite model is handed the
generator's true kernel while the direct model must estimate one from a handful
of points.

\paragraph{What is not established.} Nothing here has been tested on the real
benchmarks. Every inner function is drawn from the surrogate's own prior, and
every outer map but the multi-objective one is scalar. Whether this mechanism
explains the spread of gains in the benchmark table --- the question that
started this --- remains open, and Section~\ref{sec:mobo} is the closest we
come to it.
